\section{The adjugate method}

The Gauss--Jordan method is often the most practical method for finding inverse matrices by hand. There is also another method, built from determinants and cofactors. This method is called the adjugate method.

The adjugate method is important for two reasons. First, it gives an explicit formula for the inverse of a matrix. Second, it shows a deep connection between determinants, cofactors, and invertibility.

\subsection{The cofactor matrix and the adjugate matrix}

Let \(A=(a_{ij})\) be an \(n\times n\) matrix. For each entry \(a_{ij}\), we can compute its cofactor
\[
C_{ij}=(-1)^{i+j}M_{ij},
\]
where \(M_{ij}\) is the minor obtained by deleting row \(i\) and column \(j\).

The matrix of cofactors is
\[
C=\begin{pmatrix}
C_{11}&C_{12}&\cdots&C_{1n}\\
C_{21}&C_{22}&\cdots&C_{2n}\\
\vdots&\vdots&\ddots&\vdots\\
C_{n1}&C_{n2}&\cdots&C_{nn}
\end{pmatrix}.
\]
The adjugate matrix is the transpose of the cofactor matrix.

\begin{definitionbox}
The adjugate of a square matrix \(A\), denoted by \(\operatorname{adj}(A)\), is defined as
\[
\operatorname{adj}(A)=C^{\top},
\]
where \(C\) is the cofactor matrix of \(A\).
\end{definitionbox}

The transpose is essential. The adjugate is not simply the cofactor matrix; it is the cofactor matrix with rows and columns exchanged.

\subsection{The inverse formula}

The central formula is the following.

\begin{theorembox}
If \(A\) is a square matrix and \(\det A\ne0\), then
\[
A^{-1}=\frac{1}{\det A}\operatorname{adj}(A).
\]
\end{theorembox}

This formula should be read slowly. It says that the inverse exists only when \(\det A\ne0\). It also says that the inverse can be built from cofactors. Thus the same objects used to compute determinants also describe the inverse matrix.

A compact way to express the reason behind the formula is
\[
A\operatorname{adj}(A)=(\det A)I.
\]
If \(\det A\ne0\), we divide by \(\det A\) and obtain
\[
A\left(\frac{1}{\det A}\operatorname{adj}(A)\right)=I.
\]
So
\[
\frac{1}{\det A}\operatorname{adj}(A)
\]
is the inverse matrix.

\subsection{A \(3\times3\) example}

Let
\[
A=\begin{pmatrix}
1&0&2\\
-1&3&1\\
0&2&1
\end{pmatrix}.
\]
First compute the determinant. Expanding along the first row,
\[
\det A
=1\det\begin{pmatrix}3&1\\2&1\end{pmatrix}
+2\det\begin{pmatrix}-1&3\\0&2\end{pmatrix}.
\]
The first minor gives
\[
3\cdot1-1\cdot2=1,
\]
and the second minor gives
\[
(-1)\cdot2-3\cdot0=-2.
\]
Therefore
\[
\det A=1+2(-2)=-3.
\]
Since \(\det A\ne0\), the inverse exists.

Now compute the cofactors. We list them carefully:
\[
C_{11}=\det\begin{pmatrix}3&1\\2&1\end{pmatrix}=1,
\]
\[
C_{12}=-\det\begin{pmatrix}-1&1\\0&1\end{pmatrix}=1,
\]
\[
C_{13}=\det\begin{pmatrix}-1&3\\0&2\end{pmatrix}=-2,
\]
\[
C_{21}=-\det\begin{pmatrix}0&2\\2&1\end{pmatrix}=4,
\]
\[
C_{22}=\det\begin{pmatrix}1&2\\0&1\end{pmatrix}=1,
\]
\[
C_{23}=-\det\begin{pmatrix}1&0\\0&2\end{pmatrix}=-2,
\]
\[
C_{31}=\det\begin{pmatrix}0&2\\3&1\end{pmatrix}=-6,
\]
\[
C_{32}=-\det\begin{pmatrix}1&2\\-1&1\end{pmatrix}=-3,
\]
\[
C_{33}=\det\begin{pmatrix}1&0\\-1&3\end{pmatrix}=3.
\]
Thus the cofactor matrix is
\[
C=\begin{pmatrix}
1&1&-2\\
4&1&-2\\
-6&-3&3
\end{pmatrix}.
\]
Therefore
\[
\operatorname{adj}(A)=C^{\top}
=\begin{pmatrix}
1&4&-6\\
1&1&-3\\
-2&-2&3
\end{pmatrix}.
\]
Finally,
\[
A^{-1}=\frac{1}{-3}\begin{pmatrix}
1&4&-6\\
1&1&-3\\
-2&-2&3
\end{pmatrix}.
\]
So
\[
A^{-1}=\begin{pmatrix}
-\frac13&-\frac43&2\\
-\frac13&-\frac13&1\\
\frac23&\frac23&-1
\end{pmatrix}.
\]

This method is longer than Gauss--Jordan for this example, but it shows exactly how determinants and cofactors build the inverse.

\subsection{When should this method be used?}

For hand calculations, the adjugate method is convenient for \(2\times2\) matrices and sometimes manageable for \(3\times3\) matrices. For larger matrices it becomes computationally heavy because many cofactors must be computed.

Its conceptual value is very high: it explains why the determinant appears in inverse formulas and why zero determinant prevents invertibility.

\begin{summarybox}
The adjugate method:
\begin{enumerate}
\item Compute \(\det A\).
\item If \(\det A=0\), the inverse does not exist.
\item Compute all cofactors \(C_{ij}\).
\item Form the cofactor matrix \(C\).
\item Transpose it to obtain \(\operatorname{adj}(A)\).
\item Use \(A^{-1}=\frac{1}{\det A}\operatorname{adj}(A)\).
\end{enumerate}
\end{summarybox}
